\chapter{Introduction}
\label{cap:cap1}

The widespread reliance on manual food logging and static meal templates in 
contemporary nutrition software leaves a significant gap between user needs and 
the capabilities of current tools. To address this limitation, this thesis 
presents FitTrack, an automated nutrition management platform that integrates an 
optimization engine driven by Linear Programming with an MLP classifier to select 
and scale candidate meals based on physiological user profiles. The following 
sections detail the context that motivated this work, the objectives pursued, the 
research methodology applied, and the general structure of the thesis.

\section{Context and Importance of the Topic}
Poor dietary habits are a primary cause of global health issues like obesity, 
type 2 diabetes, and cardiovascular diseases. While growing health awareness has 
made digital nutrition tools incredibly popular—with millions of users relying 
on apps like MyFitnessPal or Cronometer—these platforms have a fundamental flaw. 
They are purely passive tracking diaries that leave all the hard work and 
nutritional decision-making to the user. 

As a result, available meal plans are either manually built or copied from rigid, 
static templates that do not adapt to a person's specific body metrics. 
To hit precise caloric and macronutrient targets, users have to change ingredient
 quantities through tedious, time-consuming trial and error. This creates a clear 
 need for smarter, automated platforms. There is a great opportunity for systems 
 that can generate precise, personalized meal plans with minimal user effort by 
 combining simple machine learning pipelines with reliable constraints optimization.

\section{General and Specific Objectives}

The primary goal of this project is to shift nutrition software from a passive 
tracking model to an automated, intelligent system. To achieve this, the objectives 
are divided into general and specific targets.

\subsection{General Objective}

The main objective of this thesis is to develop FitTrack, an automated nutrition 
management platform that eliminates the need for manual macro-nutrient guessing by 
instantly generating personalized, mathematically precise meal plans tailored to 
an individual's unique physical profile.

\subsection{Specific Objectives}

\begin{itemize}
    \item Design and implement a secure user profile system that collects and 
    processes essential physiological metrics, including age, sex, weight, and 
    activity level.
    \item Train and integrate an MLP classifier capable of evaluating user profiles 
    to select the most appropriate baseline options from a predefined batch of 
    candidate meals.
    \item Formulate and deploy a constraints-based optimization engine using Linear 
    Programming to dynamically scale ingredient weights down to the gram, ensuring 
    strict compliance with target caloric and macro-nutrient boundaries.
    \item Build a responsive and intuitive user interface that allows seamless 
    navigation, profile updates, and clear visualization of the generated meal 
    plans without administrative overhead.
    \item Evaluate the computational efficiency and mathematical accuracy of the 
    optimization solver across various target fitness scenarios.
\end{itemize}

\section{Research Methodology}

The research methodology applied in this thesis follows a structured, 
multi-phase engineering approach designed to combine machine learning 
classification with mathematical optimization.

The first phase focused on data modeling and preparation. A structured 
dataset containing a predefined batch of candidate meals was curated, 
with each meal mapped to specific nutritional baselines and ingredient 
proportions. Alongside this, user profile structures were designed to 
process essential physiological metrics such as age, sex, weight, and 
activity levels.

The second phase involved the implementation of the dual-stage 
calculation engine. The process begins with a predictive classification 
stage, where an MLP classifier is trained and deployed to analyze the 
user's physiological data and select the most appropriate baseline 
option from the meal dataset. Once the base meal is selected, 
the pipeline transitions to a deterministic optimization stage. Here, 
a constraints-based Linear Programming solver takes the target caloric 
and macro-nutrient boundaries calculated for that user and 
mathematically adjusts the individual ingredient weights down to the 
exact gram.

The final phase of the methodology centers on experimental evaluation. 
The platform's performance was assessed by running various simulated 
user profiles through the architecture to measure the solver's 
computational efficiency, verifying that it reliably satisfies all 
mathematical constraints within seconds without losing structural 
ingredient integrity.

\section{General Structure of the Thesis}

This thesis is organized into five main chapters, systematically moving 
from theoretical foundations to practical implementation and 
evaluation.

\textbf{Chapter 1} introduces the core problem, discussing the 
limitations of passive tracking tools and outlining the general and 
specific objectives of the FitTrack platform.

\textbf{Chapter 2} establishes the theoretical framework and reviews 
existing literature. It explores the fundamentals of personalized 
nutrition, the mechanics of Multi-Layer Perceptron networks, and the 
mathematical principles behind Linear Programming and constraints 
optimization.

\textbf{Chapter 3} details the architecture and system design of the 
application. This section covers the data modeling pipeline, 
the integration between the MLP classifier and the Linear Programming 
solver, and the overall software infrastructure.

\textbf{Chapter 4} focuses on the practical implementation, experimental 
results, and performance analysis. It demonstrates how the system 
scales ingredient weights under different fitness scenarios and 
evaluates the computational efficiency of the solver.

\textbf{Chapter 5} concludes the thesis by summarizing the key 
achievements of the project, discussing its limitations, and identifying 
potential directions for future development.